\section{Conclusion}
\label{sec:conclusion}
\vspace{-1mm}

We introduced \ours, a framework that bridges LLM emotion representations and EEG through soft-label knowledge distillation. By extracting emotion vectors from instruction-tuned LLMs and using their similarity structure as soft supervision targets, we achieve state-of-the-art results on FACED ($\mathbf{0.628}$ balanced accuracy, $+11.2\%$ over REVE) and significant improvements on SEED-V, with cross-subject generalization confirmed at $p < 0.001$.

We validated the framework along three orthogonal axes. (i) \textbf{Cross-backbone:} KD significantly improves both CBraMod ($+5.2\%$, $p < 10^{-4}$) and LaBraM ($+2.6\%$, $p = 0.025$), confirming architecture-agnostic benefit. (ii) \textbf{Cross-dataset:} the KD effect scales with token count, reaching $+1.85\%$ on LaBraM with 5\,s SEED-V windows ($p = 0.037$). (iii) \textbf{Cross-LLM-family:} 11 instruction-tuned models from Qwen2.5 and Gemma-4 produce nearly identical emotion geometries (within-family $\rho \approx 0.89$, cross-family $\rho = 0.82$), with $100\%$ binary cluster purity for positive vs.\ negative valence across all sufficiently capable instruction-tuned LLMs.

A central representational finding is that KD increases linear CKA between EEG features and \emph{every} LLM in our 11-model suite, including LLMs from a different family from the training teacher, with the highest post-KD alignment achieved by Gemma-4 31B (despite training on Qwen-14B). This argues that KD absorbs a universal emotion structure rather than overfitting to one specific teacher. We additionally find that freezing the pretrained backbone and training lightweight adapters matches full fine-tuning on FACED with $20\times$ fewer trainable backbone parameters, suggesting that pretrained EEG features are already rich and the critical factor is the quality of the supervision signal, and that temporal jitter is the single critical augmentation for EEG emotion classification.

Beyond classification, we demonstrate that KD shapes the \emph{geometry} of the learned EEG feature space: cluster separation improves by 38\% (Davies-Bouldin), the number of valence-correlated principal components doubles (4 vs.\ 2), and contrastive activation addition can shift emotion predictions for 92\% of samples using data-driven steering vectors. This structured manifold translates to a practical benefit: in cross-subject few-shot calibration, a KD-trained model outperforms baseline by $+5.7\%$ with just 1 labeled sample per class, reducing the calibration burden for new BCI users. We also show that mixup augmentation provides a significant additional gain in the cross-subject setting ($+1.5\%$, $p = 0.042$), stacking on top of the KD framework.

\noindent\textbf{Future Work.} The structured emotion manifold opens directions for EEG-based emotion editing and generation, where data-driven steering vectors could be used for neurofeedback-based emotion regulation. Extending \ours\ to additional EEG foundation models and emotion datasets would further validate its generality, and online adaptation of emotion vectors during training is a promising direction for personalized emotion recognition.

\noindent{\textbf{Data Availability Statement.}}
Our code, pretrained adapter weights, LLM emotion vectors for both FACED and SEED-V (across all 11 LLM models we tested), and the SEED-V 5\,s preprocessing pipeline will be made publicly available upon publication. The code will be released on GitHub.
